\chapter{Arabic-to-Latin logical transmission: texts, routes and reception}
\label{sec:p13-arabic-latin-logical-transmission}

A doctrine's appearance in two languages is not by itself evidence
that one author read the other. Transmission may involve a direct
translation, an intermediary language, incorporated excerpts,
a commentary's lemmata, or a report whose source is uncertain.
Charles Burnett's \emph{Arabic into Latin: the reception of Arabic
philosophy into Western Europe} supplies this source-bound account,
including its translation table and qualifications
\cite{Burnett2005ArabicLatinSupplement}. It is a different chapter
from his \emph{The Translation of Arabic Works on Logic into Latin
in the Middle Ages and the Renaissance} in the 2004 \emph{Handbook
of the History of Logic} \cite{Burnett2004ArabicLatinHHL}; the
two presentations are not credited as identical.

\section{What can pass through a translation}
\label{sec:p13-arabic-latin-objects}

Burnett distinguishes Arabic versions of Greek philosophical works,
questions or summaries, systematic treatises, commentaries and
doxographies. Each can carry logical material without having the
same relation to its attributed author. Avicenna's \emph{al-Shifa}
is a systematic work, not simply a commentary on Aristotle.
An Averroist long commentary instead embeds its commented text
as lemmata. A Latin reader can thereby encounter a newly translated
Aristotelian passage as well as its interpretation. The Latin title
\emph{Sufficientia} for Avicenna's work is itself a transformation
of the Arabic title, not evidence that its literal meaning was preserved.
\cite[pp.~370--371, 374--377, 387, note~6]{Burnett2005ArabicLatinSupplement}

The early natural-philosophical setting matters without becoming
the history of all logic. Constantine the African's translations
in southern Italy connected medicine and philosophy; Arabic
texts also reached scholars through psychology and astronomy.
A philosophical proposition can travel inside a work whose main
subject is not philosophy. Conversely, a work's presence in a
library does not show that every doctrine in it was adopted
by a particular reader.
\cite[pp.~371--372, 383--385]{Burnett2005ArabicLatinSupplement}

\section{Two Toledan projects, with overlap but without identity}
\label{sec:p13-arabic-latin-toledo}

Burnett explains Toledo's mid-twelfth-century prominence through
an Arabic-speaking environment, manuscript resources, Latin clergy
and Jewish scholars. The possible contribution of the displaced
Saragossan ruler's library is qualified, not a demonstrated provenance
for every text. His account distinguishes a project associated
with Gerard of Cremona from one associated with Dominicus
Gundisalvi, also called Gundissalinus.
\cite[pp.~372--374]{Burnett2005ArabicLatinSupplement}

Gerard's students compiled a list after his death in 1187.
Its subject arrangement supports Burnett's interpretation of a
curricular project: the traditional liberal arts are followed by
philosophy and medicine. Al-Farabi's \emph{Classification of the
Sciences}, which Gerard translated, supplies a relevant model
of the fields and their texts. Gerard's \emph{Posterior Analytics}
is listed as a book of demonstration, making standards for
scientific argument part of the translated curriculum. But the
students' list is not a date order of Gerard's translations.
Burnett explicitly says that none of those translations is dated.
\cite[pp.~373--375, 387, notes~10--11, 391]{Burnett2005ArabicLatinSupplement}

Gundisalvi's project gives particular prominence to Avicennian
works on soul and metaphysics, together with other Arabic texts.
Burnett considers the identification of Avendauth with Abraham
ibn Daud very likely, not unqualifiedly certain. The collaboration
on Avicenna's \emph{On the Soul} differs from the unknown Latinist
associated with Avendauth's logical sample passages. Both Toledan
projects could use al-Farabi, but different versions of some works
survive. Duplicate translations of \emph{Classification of the
Sciences}, and the Latin titles \emph{De intellectu} and
\emph{De ratione} for versions of al-Kindi's work on intellect,
show why common subject matter does not collapse the projects
into one translator or one text.
\cite[pp.~376--381, 393--395]{Burnett2005ArabicLatinSupplement}

\begin{hbexample}{def:p13-arabic-latin-list-inference}{Attribution does not supply an internal chronology}
Suppose an authenticated posthumous list attributes translations
of demonstration, metaphysics and physics to a translator and
places them in that order. It supports attribution and can support
an argument about its compiler's classification. Without dates
or further evidence it does not establish that the translator
completed metaphysics before physics. Burnett's warning about
Gerard's list is this distinction between documentary order and
the chronology of the activity recorded. The hypothetical
formulation makes the inference explicit; it is not an additional
discovered manuscript.
\end{hbexample}

\section{The logical corpus was selectively translated}
\label{sec:p13-arabic-latin-selected-logic}

The table identifies Aristotle's \emph{Posterior Analytics}
with Gerard of Cremona, and records Themistius' commentary and
an al-Farabi syllogistic work in Gerard's list. The last is
unidentified in Latin in Burnett's account. For al-Farabi's
\emph{On Interpretation} treatment it records abbreviated excerpts.
Citation of an al-Farabi work by Albert the Great is different
evidence from an identified complete Latin translation. These
records cannot all be promoted to complete surviving translations.
\cite[pp.~391--394, 402, notes~28--32]{Burnett2005ArabicLatinSupplement}

The Avicennian selection is exact rather than all-or-nothing.
Selected \emph{Isagoge} chapters translated by Avendauth with
an unknown Latinist differ from further chapters assigned to
unknown twelfth-century translators, expressly not Gundisalvi.
A section of Avicenna's demonstrative work is incorporated in
Gundisalvi's \emph{De divisione philosophiae}; excerpts from
Avicenna's rhetoric occur within Hermann the German's Aristotle
\emph{Rhetoric}. This does not establish that Avicenna's entire
modal syllogistic entered Latin in the same form.
\cite[pp.~394--395, 403, notes~40--41]{Burnett2005ArabicLatinSupplement}

Street explains why such selection is consequential. Latin
translators could favor the closer Aristotelian exegesis of
al-Farabi and Averroes while later eastern Arabic logic's central
problems were organized around Avicenna. Popularity in Latin
is not a measure of an author's relative importance in Arabic
logical traditions.
\cite[pp.~248--249]{Street2005LogicSupplement}

\begin{hbexample}{def:p13-arabic-latin-modal-route-case}{A similar distinction does not identify its route}
Street compares a Farabian restriction on necessary conversion
with Robert Kilwardby's essential/accidental necessity distinction.
The literate-human case explains their common problem: a necessary
connection in one direction does not make literacy necessary of
some human in the converse. But Street expressly says that
Kilwardby's source cannot be identified. A correspondence of
arguments therefore does not demonstrate Kilwardby's direct
reading of al-Farabi or Averroes. Transmission needs a textual
route or stronger documentary evidence than a shared distinction.
\cite[pp.~257--258]{Street2005LogicSupplement}
\end{hbexample}

\section{Excerpts, lemmata and indirect thirteenth-century knowledge}
\label{sec:p13-arabic-latin-thirteenth-century}

Michael Scot's translations and Averroes' long commentaries
allowed Latin readers to compare Aristotle in different versions.
Hermann the German's Toledo work between 1240 and 1256 also
combined an Aristotelian text with Arabic commentators. The
table identifies his 1256 version of Averroes' middle commentary
on the \emph{Poetics}; it is not an otherwise identical translation
of Aristotle's own work. A composite text can preserve several
authorities without completely translating each contributing source.
\cite[pp.~374--375, 380--382, 397]{Burnett2005ArabicLatinSupplement}

Some authors read Arabic directly. Burnett describes Ramon
Marti's citations from Arabic philosophical works in the
\emph{Pugio fidei} and Arabic knowledge in works associated
with Gonzalo Perez Gudiel and his circle. Albert's knowledge
of al-Farabi passages can likewise exceed the corpus of complete
Latin versions now identified. Burnett's qualified explanations
of this diffusion do not license invention of an otherwise
unattested full translation for an unexplained reference.
\cite[pp.~381--383]{Burnett2005ArabicLatinSupplement}

\section{Hebrew intermediaries and Renaissance printed witnesses}
\label{sec:p13-arabic-latin-hebrew-renaissance}

Later reception includes Arabic-to-Hebrew-to-Latin routes,
marked separately in Burnett's table. Renaissance knowledge
of Hebrew, Jewish scholarly collaboration and available Hebrew
versions help explain that intermediary. Linguistic proximity
was a historical rationale, not a proof of unchanged meaning.
The eleven-volume Giunta edition of Aristotle and Averroes,
published at Venice in 1550--1552, belongs to a history of
revision and new translation, not one unchanged medieval version.
\cite[pp.~385--386, 391]{Burnett2005ArabicLatinSupplement}

Logical entries include Averroes' epitomes translated through
Hebrew by Abraham de Balmes in 1523; Burana's \emph{Prior
Analytics} version in 1524; William of Luna's middle commentaries,
with Burnett's question marks preserved for some attributions;
and further Hebrew-mediated versions by Jacob Mantino.
\emph{Questions on Logic} is listed with Elias del Medigo in
1497 and Abraham de Balmes in 1523. These are distinct works
and versions. Dates printed for Renaissance editions are first
publication dates in the table, not necessarily translation dates.
\cite[pp.~391, 397--399]{Burnett2005ArabicLatinSupplement}

Burnett identifies witnesses where appropriate: Vatican Latin
12055 for a Hebrew-mediated al-Farabi intellect text; Vatican
Latin 2186 for a fragmentary al-Farabi-related work; and Vatican
Reginensis Latin 1268, folios~72r--73r, for al-Farabi material
incorporated in a Euclidean commentary. These are reported
identifications, not local manuscript inspections. A passage's
host can matter as much as its conventional standalone title.
\cite[pp.~393--394, 402--403, notes~29 and 33]{Burnett2005ArabicLatinSupplement}

The table ends with the collection \emph{Epistolae seu Quesita
Logica Diversorum Doctorum Arabum precipue Averroys}, including
entries whose Arabic originals are not known and whose last
three authors were not identified in Burnett's account. Printed
association with Arabic logic cannot provide secure biographies
or affiliations for those authors. An asterisk can mark a lost
text or unidentified item; it does not erase positive evidence
for an earlier translation.
\cite[pp.~400, 404, note~68]{Burnett2005ArabicLatinSupplement}

\begin{hbmethod}{def:p13-arabic-latin-transmission-method}{Evaluate a claimed logical transmission}
Identify the Arabic work and passage, its genre and selected
version, and the Latin witness or explicit citation. Separate
translator attribution from survival, and direct Arabic
translation from Hebrew mediation. Record uncertain names,
composite excerpts and different Latin versions rather than
selecting one to make a simple lineage. Date composition,
translation, surviving copying and printed publication separately.
Compare quantifier scope, existence, modal condition and technical
meaning before calling two inferences the same rule. A
correspondence can motivate investigation; it does not alone
prove access, borrowing or unchanged doctrinal reception.
\end{hbmethod}

\begin{proofboundary}{def:p13-arabic-latin-source-boundary}{The extent of the evidence}
These particulars come from Burnett's complete 2005 companion
chapter and its qualified table, supplemented by Street's
identified modal comparison. Burnett's distinct ten-page 2004
logic-translation argument is treated in
Chapter~\ref{sec:p13-burnett-arabic-latin-2004}. A translation's existence does not
establish a theorem about the equivalence of its terminology.
Such a comparison needs the source and target passages;
a bibliographical route and a modern valid argument are
different evidence.
\end{proofboundary}
